\section{Całka Riemanna--Stieltjesa}
\begin{definicja}
Niech $f,\alpha:[a,b]\to\R$. Piszemy $f\in\mathcal R(\alpha)$, jeśli istnieje $I\in\R$ takie, że dla każdego $\varepsilon>0$ istnieje $\delta>0$ i dla każdego podziału znaczonego $\mathcal P$ o $\mesh P<\delta$ zachodzi
\[
\left|\sum_{i=1}^n f(t_i)(\alpha(x_i)-\alpha(x_{i-1}))-I\right|<\varepsilon.
\]
Wtedy $I=\int_a^b f\,\mathrm d\alpha$ jest całką Riemanna--Stieltjesa. Przy $\alpha(x)=x$ otrzymujemy całkę Riemanna.
\end{definicja}
\begin{uwaga}[Rozszerzenie]
Dla niemalejącej $\alpha$ i ograniczonej $f$ można definiować sumy dolne i górne Stieltjesa z przyrostami $\alpha$. Jeżeli $\alpha$ ma skok, nawet ciągła w sensie jednostronnym funkcja $f$ może wymagać zgodności znaczników w otoczeniu tego skoku; zwykła całka Riemanna--Stieltjesa może wówczas nie istnieć.
\end{uwaga}
Dla ograniczonej funkcji $f$ i niepustego zbioru $J\subset[a,b]$ oznaczamy $\osc_J f=\sup_J f-\inf_J f$. Lokalna oscylacja to $\osc_x f=\inf_{\rho>0}\osc_{(x-\rho,x+\rho)\cap[a,b]}f$. Jest równa zeru dokładnie w punktach ciągłości względnej $f$: jeśli oscylacja na otoczeniu jest mała, to w szczególności wszystkie wartości w nim są bliskie $f(x)$; odwrotnie ciągłość daje małą różnicę każdej pary wartości. Dla $r>0$ zbiór $D_r=\{x:\osc_x f\ge r\}$ jest domknięty, ponieważ gdy $\osc_x f<r$, odpowiednio małe otoczenia punktów pobliskich zawierają się w otoczeniu $x$ o oscylacji mniejszej od $r$. W szczególności $D(f)=\bigcup_{n\ge1}D_{1/n}$ jest zbiorem borelowskim.

\begin{twierdzenie}[Narzędzia]\label{it:thm:rs-tools}
Niech $V_a^b(\alpha)=\sup_P\sum_i|\alpha(x_i)-\alpha(x_{i-1})|$.
\begin{enumerate}[label=(\alph*)]
\item Jeśli $f$ jest ciągła i $\alpha$ ma wahanie ograniczone, to $f\in\mathcal R(\alpha)$ i $\bigl|\int_a^b f\,\mathrm d\alpha\bigr|\le\|f\|_\infty V_a^b(\alpha)$.
\item Jeśli $f,\alpha$ mają wahanie ograniczone i nie mają wspólnych punktów nieciągłości, to obie całki Stieltjesa istnieją oraz
$\int_a^b f\,\mathrm d\alpha+\int_a^b\alpha\,\mathrm df=f(b)\alpha(b)-f(a)\alpha(a)$.
\item Jeśli $\alpha\in C^1[a,b]$ i $f\in\mathcal R[a,b]$, to $\int_a^b f\,\mathrm d\alpha=\int_a^b f(x)\alpha'(x)\,\mathrm dx$.
\item Całka jest liniowa w każdej zmiennej, jeśli odpowiednie całki istnieją. Dla niemalejącej $\alpha$ jest dodatnia i monotoniczna względem $f$.
\end{enumerate}
\end{twierdzenie}
\begin{proof}
Dla (a) rozbijamy każdą komórkę jednego podziału według punktów drugiego. Różnica sum przy zagęszczeniu jest ograniczona przez $\omega_f(\mesh P)V_a^b(\alpha)$, gdzie $\omega_f(h)=\sup_{|s-t|\le h}|f(s)-f(t)|$. Porównując dwa podziały z ich wspólnym zagęszczeniem i korzystając z $\omega_f(h)\to0$, dostajemy kryterium Cauchy'ego. Oszacowanie całki wynika z oszacowania każdej sumy.

Dla (b) udowodnimy najpierw istnienie całki. Oznaczmy $V(x)=\Var_{[a,x]}\alpha$ oraz $T=V(b)$. Z addytywności wariacji mamy $V(v)-V(u)=\Var_{[u,v]}\alpha$. Funkcja $V$ jest ciągła w każdym punkcie ciągłości $\alpha$. Istotnie, dla $x<b$ wybieramy podział $[x,b]$, którego suma wariacyjna różni się od $\Var_{[x,b]}\alpha$ o mniej niż $\eta>0$. Dla $y>x$ przed pierwszym punktem tego podziału, zastąpienie jego początku $x$ przez $y$ zmniejsza sumę o co najwyżej $|\alpha(y)-\alpha(x)|$. Z addytywności wariacji wynika
\[
 V(y)-V(x)\le\eta+|\alpha(y)-\alpha(x)|.
\]
Przechodząc do granicy przy $y\downarrow x$, a następnie $\eta\downarrow0$, uzyskujemy ciągłość prawostronną $V$ w punkcie ciągłości $\alpha$. Analogiczny argument dla podziału $[a,x]$ daje ciągłość lewostronną. Na końcach stosujemy tylko odpowiedni kierunek.

Niech $r>0$. Zbiór punktów $D_r$, w których lokalna oscylacja $f$ jest co najmniej $r$, jest skończony. Dla dowolnej skończonej rodziny różnych punktów $D_r$ wybieramy rozłączne względne otoczenia przedziałowe, w każdym z nich dwa punkty o różnicy wartości $f$ większej od $r/2$. Suma tych różnic nie przekracza $\Var_{[a,b]}f$, więc liczba wybranych punktów nie przekracza $2\Var_{[a,b]}f/r$. To daje skończoność $D_r$.

Ponieważ $r>0$, punkty $D_r$ są punktami nieciągłości $f$, a więc punktami ciągłości $\alpha$ i $V$. Wybieramy dla każdego z nich względne otwarte otoczenie przedziałowe $O_x$ o domknięciu tak małym, że suma wariacji $\alpha$ na domknięciach wszystkich $O_x$ jest mniejsza od $\eta$. Dla każdego punktu spoza $D_r$ wybieramy względne otwarte otoczenie, na którym oscylacja $f$ jest mniejsza od $r$. Otrzymujemy otwarte pokrycie zwartego przedziału. Z lematu o liczbie Lebesgue'a istnieje $h>0$ takie, że każda komórka o długości mniejszej od $h$ leży w jednym z tych otoczeń. Komórki przypisane do $O_x$ mają łączną wariację $\alpha$ nie większą niż suma wariacji na domknięciach $O_x$, ponieważ mają rozłączne wnętrza i wariacja jest addytywna. Pozostałe komórki mają oscylację $f$ mniejszą od $r$. Zatem dla każdego takiego podziału $P$,
\[
 \sum_{I\in P}\osc_I f\,\Var_I\alpha
 \le 2\|f\|_\infty\eta+rT.
\]
Wybór najpierw małego $r$, a następnie małego $\eta$, pozwala uczynić prawą stronę dowolnie małą (także gdy $T=0$ lub $f=0$).

Dla dowolnego rozdrobnienia $P'$ podziału $P$ różnica sum Stieltjesa, niezależnie od znaczników, ma moduł nie większy niż lewa strona ostatniej nierówności: rozpisujemy przyrost $\alpha$ na każdej starej komórce jako sumę przyrostów na nowych komórkach i stosujemy nierówność trójkąta. Porównanie dwóch dostatecznie drobnych podziałów ze wspólnym rozdrobnieniem daje kryterium Cauchy'ego dla sum. Istnienie wspólnej granicy wynika z zupełności $\R$: wybieramy ciąg podziałów o średnicach dążących do zera, ich sumy tworzą ciąg Cauchy'ego, a każdą dostatecznie drobną sumę porównujemy z jego granicą. Tak otrzymujemy $\int f\,\mathrm d\alpha$. Zamiana ról $f,\alpha$ daje istnienie $\int\alpha\,\mathrm df$.

Dla dowolnego podziału zachodzi dokładna tożsamość
\[
 \sum_i f(x_{i-1})(\alpha(x_i)-\alpha(x_{i-1}))
 +\sum_i\alpha(x_i)(f(x_i)-f(x_{i-1}))
 =f(b)\alpha(b)-f(a)\alpha(a).
\]
Wynika ona z rozpisania każdego składnika jako różnicy kolejnych iloczynów i teleskopowania. Przechodząc do granicy przy średnicy dążącej do zera, otrzymujemy wzór całkowania przez części. W (c) twierdzenie Lagrange'a daje $\alpha(x_i)-\alpha(x_{i-1})=\alpha'(\xi_i)(x_i-x_{i-1})$; z jednostajnej ciągłości $\alpha'$ wynika, że możemy zastąpić $\xi_i$ przez $t_i$. Ostatnia teza wynika bezpośrednio z definicji.
\end{proof}
\begin{twierdzenie}[Kryterium oscylacyjne]\label{it:thm:rs-oscillation}
Niech $\alpha:[a,b]\to\R$ będzie niemalejąca, prawostronnie ciągła, a $f:[a,b]\to\R$ ograniczona. Niech $\mu_\alpha$ będzie uzupełnioną miarą Lebesgue'a--Stieltjesa określoną przez $\mu_\alpha((s,t])=\alpha(t)-\alpha(s)$ dla $a\le s<t\le b$, przy $\mu_\alpha(\{a\})=0$. Następujące warunki są równoważne:
\begin{enumerate}[label=(\roman*)]
\item $f\in\mathcal R(\alpha)$;
\item $\mu_\alpha(D(f))=0$;
\item dla każdego $\varepsilon>0$ istnieje $\delta>0$ takie, że dla każdego podziału $P$ o $\mesh P<\delta$ zachodzi $\sum_i\osc_{[x_{i-1},x_i]}f\,(\alpha(x_i)-\alpha(x_{i-1}))<\varepsilon$.
\end{enumerate}
Wówczas $\int_a^b f\,\mathrm d\alpha=\int_{(a,b]}f\,\mathrm d\mu_\alpha$.
\end{twierdzenie}
\begin{proof}
Dla niemalejącej $\alpha$ suma górna minus dolna jest dokładnie sumą z (iii). Porównywanie sum dla dowolnych znaczników dowodzi (i)$\Leftrightarrow$(iii), tak samo jak w kryterium Darboux.

Niech $M=\|f\|_\infty$. Jeżeli $\mu_\alpha(D(f))=0$, to dla $r>0$ zwarty zbiór $D_r=\{x:\osc_x f\ge r\}$ można pokryć zbiorem otwartym $G$ o dowolnie małej mierze $\mu_\alpha$. Poza $G$ oscylacje $f$ są lokalnie mniejsze od $r$. Pokrycie przez $G$ i względne otwarte otoczenia punktów spoza $D_r$, na których oscylacja jest mniejsza od $r$, ma dodatnią liczbę Lebesgue'a. Przy średnicy mniejszej od tej liczby każda komórka mieści się w jednym z tych zbiorów. Przedziały przypisane do $G$ wnoszą do sumy (iii) nie więcej niż $2M\mu_\alpha(G)$, a pozostałe nie więcej niż $r(\alpha(b)-\alpha(a))$. To dowodzi (iii).

Jeśli zachodzi (iii), wybieramy podziały z coraz mniejszą średnicą i punktami podziału poza atomami $\mu_\alpha$ (można to zrobić, ponieważ atomów jest co najwyżej przeliczalnie wiele). Komórki o wnętrzach przecinających $D_r$ mają oscylację co najmniej $r$, więc ich suma miar jest arbitralnie mała. Punkty wewnętrzne leżące na granicach komórek nie są atomami, zaś obecność punktu $b$ w $D_r$ powoduje oscylację w ostatniej komórce. Przy konwencji $\mu_\alpha(\{a\})=0$ otrzymujemy $\mu_\alpha(D_r)=0$. Sumowanie po $r=1/n$ daje (ii).

Dla dowolnego ciągu podziałów o średnicach dążących do zera definiujemy funkcje schodkowe dolne i górne na komórkach półotwartych $(x_{i-1},x_i]$, używając infimum i supremum po domkniętych komórkach. W każdym punkcie ciągłości $f$ obie funkcje zbiegają do $f$. Ponieważ $D(f)$ jest zbiorem $\mu_\alpha$-zerowym, uzupełnienie miary zapewnia mierzalność $f$ (poza tym zbiorem jest granicą funkcji schodkowych). Funkcje schodkowe zbiegają zatem prawie wszędzie do $f$ i mają całki Lebesgue'a--Stieltjesa równe sumom dolnym i górnym. Zbieżność zmajoryzowana daje ostatnią równość.
\end{proof}
\begin{uwaga}[Atomy i końce przedziału]
Jeżeli $\mu_\alpha(\{c\})=\alpha(c)-\alpha(c-)>0$ dla $c\in(a,b]$, warunek z poprzedniego twierdzenia wymusza ciągłość $f$ w $c$ (przy $c=b$ w sensie lewostronnym). Konwencja $\mu_\alpha(\{a\})=0$ wynika z użycia przyrostów od $\alpha(a)$; przy innej konwencji konieczna jest osobna poprawka za atom w $a$.
\end{uwaga}
\begin{przyklad}[Część ciągła i atomy]
Na $[0,1]$ weźmy $\alpha(x)=x+\tfrac12\ind{[1/3,1]}(x)+\ind{[2/3,1]}(x)$. Dla $f(x)=x^2$ istnieje całka Stieltjesa i
\[
\int_0^1x^2\,\mathrm d\alpha=\int_0^1x^2\,\mathrm dx+\frac12\left(\frac13\right)^2+\left(\frac23\right)^2=\frac13+\frac1{18}+\frac49=\frac56.
\]
Dla $\alpha(x)=\lfloor x\rfloor$ na $[0,n]$, $n\in\N$, oraz ciągłej $f$ otrzymujemy analogicznie $\int_0^n f\,\mathrm d\alpha=\sum_{k=1}^n f(k)$.
\end{przyklad}
\begin{przyklad}[Wspólny skok]
Niech $c\in(a,b)$ i $f=\alpha=\ind{[c,b]}$. Zwykła całka Riemanna--Stieltjesa nie istnieje: na przedziale zawierającym $c$ suma zależy od tego, czy znacznik wybrano przed czy po $c$. Natomiast pojęcie całki Kurzweila--Stieltjesa pozwoli otrzymać wartość $1$.
\end{przyklad}
\begin{przyklad}[Dlaczego nie wystarczy wybrać jednego podziału]
Niech $a<c<b$, $\alpha=\ind{[c,b]}$ i $f=\ind{(c,b]}$, przy czym $f(c)=0$. Dla podziału zawierającego $c$ każda suma dolna i górna Stieltjesa może być równa zeru: dodatni przyrost funkcji całkującej kończy się w $c$, a na $[x_{i-1},c]$ funkcja $f$ jest równa zeru. Podziały przecinające $c$ mogą jednak dawać sumy $0$ albo $1$, zależnie od znacznika. Warunek „dla wszystkich dostatecznie drobnych podziałów” w definicji Riemanna--Stieltjesa jest więc istotny; zagęszczanie jednego wybranego podziału nie wystarcza.
\end{przyklad}
\begin{przyklad}[Wpływ funkcji całkującej na całkowalność]
Na $[0,1]$ niech $f=\ind{\Q\cap(0,1/4]}$ i $\alpha=\ind{[1/2,1]}$. Funkcja $f$ nie jest całkowalna Riemanna względem długości, gdyż jest nieciągła na całym $(0,1/4)$. Jest jednak ciągła i równa zeru w otoczeniu jedynego atomu $\mu_\alpha$, zatem $\int_0^1 f\,\mathrm d\alpha=0$ w sensie Riemanna--Stieltjesa. Z kolei funkcja z przykładu 3.7 jest Riemanna całkowalna względem długości, ale nie względem wskazanego $\alpha$. Dla różnych funkcji całkujących klasy całkowalności nie muszą być porównywalne przez inkluzję.
\end{przyklad}

% ============================================================
% UZUPELNIENIE MATEMATYCZNE: ROZDZIAL 3
% ============================================================
\subsection*{Rachunek praktyczny: skoki i kontrola błędu}
\begin{twierdzenie}[Jeden skok w zwykłej całce Stieltjesa]\label{tw:rs-jeden-skok}
Niech $a<c<b$, $A\ne0$, $\alpha(x)=A\ind{[c,b]}(x)$ dla $x\in[a,b]$, a $f:[a,b]\to\R$ będzie ograniczona. Całka Riemanna--Stieltjesa $\int_a^b f\,\mathrm d\alpha$ istnieje wtedy i tylko wtedy, gdy $f$ jest ciągła w $c$. Wówczas $\int_a^b f\,\mathrm d\alpha=Af(c)$. Jeśli skok jest umieszczony w $b$, warunkiem jest ciągłość lewostronna $f$ w $b$, a dla skoku umieszczonego w $a$ funkcją $\ind{[a,b]}$ integrator jest stały i całka wynosi zero.
\end{twierdzenie}
\begin{proof}
W każdej sumie Stieltjesa wszystkie przyrosty $\alpha$ są zerowe poza jedną komórką $[u,v]$, dla której $u<c\le v$. Suma wynosi więc $Af(t)$, gdzie $t\in[u,v]$. Jeśli $f$ jest ciągła w $c$, dostatecznie mała średnica podziału wymusza $|f(t)-f(c)|<\varepsilon/|A|$.

Odwrotnie, załóżmy istnienie całki i rozważmy dowolny ciąg $t_n\to c$. Dobieramy komórki $[u_n,v_n]$ o długościach dążących do zera, zawierające zarówno $c$, jak i $t_n$, przy czym $u_n<c\le v_n$. Uzupełniamy je do podziałów o średnicach dążących do zera. Sumy ze znacznikiem $t_n$ oraz ze znacznikiem $c$ muszą mieć tę samą granicę. Zatem $A(f(t_n)-f(c))\to0$. Ponieważ $(t_n)$ był dowolny, $f$ jest ciągła w $c$. Wariant dla $c=b$ ma identyczny dowód z granicą jednostronną. W punkcie $a$ integrator $A\ind{[a,b]}$ ma wszystkie przyrosty równe zeru.
\end{proof}
\begin{wniosek}[Skończona liczba atomów]\label{wn:rs-suma-skokow}
Niech $\alpha(x)=h(x)+\sum_{k=1}^m A_k\ind{[c_k,b]}(x)$ na $[a,b]$, gdzie $h\in\BV[a,b]$ jest ciągła, $a<c_1<\cdots<c_m\le b$, a $A_k\in\R$. Jeśli $f\in C[a,b]$, to
\[
\int_a^b f\,\mathrm d\alpha=\int_a^b f\,\mathrm dh+\sum_{k=1}^m A_kf(c_k).
\]
W przypadku samej części skokowej wystarczy ciągłość $f$ we wszystkich $c_k$ (jednostronna w $b$).
\end{wniosek}
\begin{proof}
Stosujemy liniowość całki i twierdzenie \ref{tw:rs-jeden-skok}. Dla ciągłego $f$ całka względem $h\in\BV$ istnieje na mocy twierdzenia z początku rozdziału.
\end{proof}
\begin{twierdzenie}[Oszacowanie błędu sumy Stieltjesa]\label{tw:rs-blad-kwadratury}
Niech $f\in C[a,b]$, $\alpha\in\BV[a,b]$, a $\mathcal P$ będzie dowolnym podziałem znaczonym o średnicy $h$. Wówczas
\[
\left|S(f,\alpha;\mathcal P)-\int_a^b f\,\mathrm d\alpha\right|
\le \omega_f(h)\Var_a^b(\alpha),
\qquad
\omega_f(h):=\sup_{\substack{x,y\in[a,b]\\|x-y|\le h}}|f(x)-f(y)|.
\]
Jeżeli $f$ jest lipschitzowska ze stałą $L$, błąd nie przekracza $Lh\Var_a^b(\alpha)$.
\end{twierdzenie}
\begin{proof}
Ustalmy podział $\mathcal P$ i zagęszczajmy każdą jego komórkę niezależnie, wybierając dowolne nowe znaczniki. Wewnątrz komórki $[x_{i-1},x_i]$ oba rodzaje znaczników są oddalone najwyżej o $h$; zatem różnica starego wkładu i sumy po jej zagęszczeniu ma moduł nie większy niż $\omega_f(h)\Var_{x_{i-1}}^{x_i}(\alpha)$. Po zsumowaniu otrzymujemy podane oszacowanie dla różnicy sum. Przy średnicy zagęszczenia dążącej do zera nowa suma dąży do całki Riemanna--Stieltjesa.
\end{proof}
\begin{przyklad}[Skończona kwadratura ze znanym błędem]
Niech $f(x)=e^x$ na $[0,1]$ i $\alpha(x)=x+2\ind{[1/2,1]}(x)$. Wówczas
$\int_0^1e^x\,\mathrm d\alpha=e-1+2\sqrt e$.
Dla dowolnego podziału znaczonego o średnicy $h$ twierdzenie \ref{tw:rs-blad-kwadratury} daje
$|S-(e-1+2\sqrt e)|\le 3eh$, gdyż $\Var_0^1(\alpha)=3$ i $\omega_f(h)\le eh$.
\end{przyklad}
\begin{uwaga}[Zgodność trzech konwencji]
W zwykłej całce Riemanna--Stieltjesa przy skoku w $c$ dowolne znaczniki wymuszają ciągłość $f$ w $c$. W całce Kurzweila--Stieltjesa cechowanie może wymusić znacznik $c$ i wówczas poprawny wkład skoku wynosi $f(c)\Delta^-\alpha(c)$, nawet jeśli $f$ nie jest ciągła. Całka Lebesgue'a--Stieltjesa wymaga mierzalności i całkowalności względem miary, lecz nie ciągłości w atomie: wkład wynosi również $f(c)\mu_\alpha(\{c\})$. Zbieżność wartości całek nie oznacza więc równości ich dziedzin.
\end{uwaga}

